% ============================================================
%  chapters/ch23-chain-of-memory.tex
% ============================================================

\chapter{Chain of Memory}

\chapterepigraph{%
  You are not the sum of your memories.\\
  You are the process that generates\\
  the next memory from the last.
}{Flyxion, \textit{Semantic Infrastructure}}

\sidenote{App.\,L §L.22\\on the chain\\of memory;\\connects to\\Appendix~M\\on composition}

Chapter~22 established that an agent's experienced world is a latent
construct — a sparse simulation maintained over time. This chapter examines
the temporal dimension of this process: how the simulation persists, how
it is updated, and what it means for an agent to have a continuous identity
through time.

The key claim is that memory is not storage but \emph{constraint
preservation}: identity over time is not the persistence of a fixed data
structure but the continuity of a process that generates each moment from
the constraints established by previous moments.

\section{Memory as Constraint Preservation}

The conventional picture of memory is archival: experiences are encoded,
stored, and later retrieved. This picture is demonstrably false in its
details — human memory is reconstructive, not reproductive; it is
influenced by subsequent experience, current mood, and retrieval context
in ways that no archive would permit.

The constraint-preservation picture is more accurate. What persists is
not a record of the experience but a \emph{modification of the constraint
structure} governing future experience.

\begin{definition}{Constraint-Preserving Memory}{constraint-memory}
  A \emph{constraint-preserving memory system} is a dynamical system
  $(X, \mathcal{C}, \tau)$ where:
  \begin{itemize}
    \item $X$ is the agent's state space,
    \item $\mathcal{C}(t)$ is the active constraint set at time $t$,
    \item $\tau : \mathcal{C}(t) \times E_t \to \mathcal{C}(t+1)$
          is the constraint update function that integrates new experience
          $E_t$ into the existing constraint structure.
  \end{itemize}
  A \emph{memory} of event $E_t$ is the modification
  $\Delta\mathcal{C} = \mathcal{C}(t+1) - \mathcal{C}(t)$ that $E_t$
  induces — the additional constraints it adds.
\end{definition}

Under this definition, forgetting is constraint relaxation: the constraint
modification induced by $E_t$ decays over time, gradually losing its
influence on the active constraint set. A vivid memory is a strong,
persistent constraint modification; a faded memory is a weak, partially
decayed one.

\begin{historynote}
  Ibn Rushd's (Averroës) commentary on Aristotle's \textit{De Memoria}
  distinguishes between the \emph{form} of a memory (its content) and
  the \emph{material disposition} (the physical substrate altered by
  experience). The constraint-preservation view maps onto this distinction:
  the constraint modification $\Delta\mathcal{C}$ is the material
  disposition; the experience $E_t$ it was derived from is the form.
  What persists is the disposition — the constraint — not a reproduction
  of the original experience.
\end{historynote}

\section{Trajectory Reconstruction}

\begin{tcolorbox}[enhanced, breakable,
  colback=RSVPlight, colframe=RSVPmid, arc=4pt, boxrule=0.6pt,
  title={\bfseries\sffamily Illustration: Memory as Stabilized Field Residue},
  coltitle=white, attach boxed title to top left={yshift=-2mm,xshift=6mm},
  boxed title style={colback=RSVPmid,sharp corners},
  top=4mm,left=4mm,right=4mm,bottom=4mm]

The HYDRA framework for cognitive field theory offers a vivid
implementation of memory-as-constraint-preservation.

When an experience is sufficiently salient — high emotional or semantic
weight — it does not merely pass through the cognitive manifold and vanish.
It deposits a \emph{stabilized field residue}: a persistent standing wave
in the scalar field $\Phi_s$ with four properties:

\begin{itemize}
  \item \textbf{Amplitude:} the intensity or salience of the memory.
        High-amplitude residues (traumatic events, moments of insight)
        persist for decades; low-amplitude residues decay in days.
  \item \textbf{Frequency:} the semantic content — the specific
        conceptual ``pitch'' that the memory encodes.
  \item \textbf{Phase:} alignment with the vector flow $\mathbf{v}_s$,
        determining associative links to neighboring memories.
  \item \textbf{Decay rate:} governed by the entropy field $S_s$.
        High-entropy residues (vague, unfocused) decay quickly;
        low-entropy residues (tightly constraining future thoughts)
        persist.
\end{itemize}

Crucially, \emph{retrieval is resonance, not lookup}. When you try to
remember something, the cognitive system broadcasts a retrieval cue — a
probe wave at a specific frequency. Standing waves in the manifold that
resonate with the probe's frequency and phase amplify; those that do not,
remain silent. The memory is reconstructed through geometric resonance,
not address-based fetch.

This explains why retrieval is context-sensitive: the echo is shaped by
the current state of the cognitive field. The same cue broadcast in
different emotional or cognitive states produces different resonances,
which is why the same event can be ``remembered'' differently at different
times. Memory was never a fixed record — it was always a reconstruction
shaped by the current geometry.

This also explains the cascade structure of memory loss in dementia:
the standing waves are not deleted but progressively lose amplitude and
phase coherence, so retrieval cues find weaker, blurrier resonances rather
than sharp, precise ones.
\end{tcolorbox}


Memory allows the agent to reconstruct not just past events but past
trajectories — the paths through state space that led to the current moment.

\begin{definition}{Trajectory Reconstruction}{traj-reconstruct-mem}
  Given the current constraint set $\mathcal{C}(t)$ and the current state
  $x(t)$, \emph{trajectory reconstruction} is the inverse inference
  problem: find the distribution over past trajectories
  $\{\gamma : [0,t] \to X\}$ consistent with $\mathcal{C}(t)$ and $x(t)$.
\end{definition}

This is the memory-as-inverse-inference picture. Remembering is not
retrieval from an archive but reconstruction from current constraints.
The constraints encode what the past \emph{must have been like} to
produce the current state — not what it actually was.

This explains why memories are malleable. Adding new constraints (learning
new facts, having new experiences, being exposed to suggestion) changes
the inverse inference problem and therefore changes what is ``remembered.''
The memory was never fixed — it was always a reconstruction.

\section{Identity Over Time}

If memory is constraint preservation rather than archival, then personal
identity over time is a different kind of thing than is usually assumed.

\begin{definition}{Process Identity}{process-identity}
  An agent has \emph{process identity} over an interval $[t_0, t_1]$
  if there exists a continuous trajectory $\gamma : [t_0, t_1] \to X$
  such that the agent's constraint set $\mathcal{C}(t)$ evolves
  continuously along $\gamma$ — each constraint state is reachable from
  the previous one by a single constraint-update step.
\end{definition}

Process identity does not require that the agent's state be the same at
$t_0$ and $t_1$ — it only requires that the states be connected by a
continuous chain of constraint updates. A person who changes their beliefs,
values, and memories radically over decades still has process identity if
each step in the change is a continuous constraint modification.

\begin{remark}
  This gives a precise answer to the Ship of Theseus problem for
  personal identity. The question is not whether any particular component
  (memory, belief, personality trait) persists unchanged, but whether the
  sequence of constraint updates connecting the earlier and later person
  is continuous. A gradual change that preserves process identity is
  ``the same person''; a discontinuous jump (severe amnesia, radical
  personality change from brain injury) may break process identity.
\end{remark}

\section{The Chain Operator}

The formal structure of memory as constraint preservation is captured
by the \emph{chain operator}: the composition of all constraint updates
from the agent's origin to the present.

\begin{definition}{Chain Operator}{chain-op}
  The \emph{chain operator} at time $t$ is
  \[
    \mathcal{C}(t) = \tau^{(t)} \circ \tau^{(t-1)} \circ \cdots
    \circ \tau^{(0)}(\mathcal{C}_0)
  \]
  where $\mathcal{C}_0$ is the initial constraint set and each
  $\tau^{(s)} = \tau(\cdot, E_s)$ is the constraint update induced
  by experience $E_s$.
\end{definition}

The chain operator makes explicit that the current constraint set is
the accumulated product of all past experiences — not stored individually
but composed into a single operator. The chain \emph{is} the memory;
individual memories are the factors.

\begin{theorem}{Composition Determines State}{composition-state}
  The agent's current state $x(t)$ is a function of $\mathcal{C}(t)$
  and the current observation $O_t$:
  $x(t) = f(\mathcal{C}(t), O_t)$.
  Therefore the agent's state is entirely determined by the chain of
  past experiences and the current observation — nothing more.
\end{theorem}

\begin{proof}
  By induction: the initial state $x(0) = f(\mathcal{C}_0, O_0)$
  is determined by initial constraints and initial observation. Each
  subsequent state $x(t+1) = f(\mathcal{C}(t+1), O_{t+1}) =
  f(\tau(\mathcal{C}(t), E_t), O_{t+1})$ is determined by the
  previous constraint set (which by induction encodes all previous
  experiences), the current experience, and the current observation.
\end{proof}

\section{Collective Memory and Semantic Infrastructure}

The chain-of-memory framework extends naturally from individual agents
to collective systems. A civilization's collective memory is the chain
of constraint updates accumulated through institutions, documents,
practices, and traditions.

\begin{definition}{Civilizational Memory Chain}{civ-memory-chain}
  A \emph{civilizational memory chain} is a chain operator
  $\mathcal{C}_{\text{civ}}(t)$ defined over a community of agents
  $(A_1, \ldots, A_n)$ with shared constraint update mechanisms:
  transmission (communication), storage (writing, recording), and
  revision (scholarship, reinterpretation).
\end{definition}

The fragility of civilizational memory corresponds to breaks in this
chain: the burning of libraries (Alexandria), the suppression of
languages, the erasure of traditions. Each break is a discontinuity
in the chain operator — the accumulated constraints of previous
generations failing to propagate to the next.

The semantic infrastructure of Chapter~25 and Appendix~M is the
engineering of civilizational memory chains: the design of transmission,
storage, and revision mechanisms that preserve constraint continuity
across generations.

\begin{exercise}
  Formalize ``forgetting'' as a constraint-relaxation operator:
  $F_\lambda(\mathcal{C}) = \{c \in \mathcal{C} : \text{strength}(c) > \lambda\}$
  where $\lambda$ is a forgetting threshold. Show that the Ebbinghaus
  forgetting curve (exponential decay of memory strength) corresponds
  to a specific choice of strength decay dynamics.
\end{exercise}

\begin{exercise}
  Two agents $A$ and $B$ begin with the same initial constraint set
  $\mathcal{C}_0$ but undergo different experience sequences.
  Define a metric $d(\mathcal{C}_A(t), \mathcal{C}_B(t))$ measuring
  how far apart their constraint sets have diverged at time $t$.
  Under what conditions does this divergence grow unboundedly?
  Under what conditions do the constraint sets reconverge?
\end{exercise}

\begin{exercise}[title={Philosophical}]
  The Buddhist concept of \emph{anatt\=a} (non-self) holds that there
  is no persistent, unchanging self — only a stream of causally connected
  mental events. Is this compatible with process identity as defined
  above? Does process identity require a ``self,'' or is it precisely the
  formal structure that makes non-self coherent?
\end{exercise}
